\chapter{Anhang}
\label{chap:anhang}


\section{Schaltpläne}
\label{sec:anhang-schaltplaene}

\begin{figure}[p]
    \centering
    \rotatebox{90}{\includegraphics[width=0.80\textheight]{images/anhang/sch_bridge.pdf}}
    \caption{Schaltplan der Bridge (Abschnitt~\ref{subsec:bridge-pcb}).}
    \label{fig:sch-bridge}
\end{figure}

\begin{figure}[p]
    \centering
    \rotatebox{90}{\includegraphics[width=0.80\textheight]{images/anhang/sch_tube.pdf}}
    \caption{Schaltplan der Leuchtröhre (Abschnitt~\ref{subsec:tube-mainboard}).}
    \label{fig:sch-tube}
\end{figure}

\begin{figure}[p]
    \centering
    \rotatebox{90}{\includegraphics[width=0.80\textheight]{images/anhang/sch_interface.pdf}}
    \caption{Schaltplan des Interface-PCB (Abschnitt~\ref{subsec:interface-pcb}).}
    \label{fig:sch-interface}
\end{figure}

\clearpage

\section{Leiterplattenlayouts}
\label{sec:anhang-layouts}

Dargestellt sind Kupferlage und Bestückungsdruck; die Unterseiten sind
spiegelbildlich, also in Durchsicht von oben, gezeichnet.

\begin{figure}[htbp]
    \centering
    \begin{subfigure}[b]{\linewidth}
        \centering
        \includegraphics[width=0.92\linewidth]{images/anhang/pcb_bridge_top.pdf}
        \caption{Oberseite (F.Cu)}
    \end{subfigure}

    \vspace{1em}

    \begin{subfigure}[b]{\linewidth}
        \centering
        \includegraphics[width=0.92\linewidth]{images/anhang/pcb_bridge_bot.pdf}
        \caption{Unterseite (B.Cu)}
    \end{subfigure}
    \caption{Leiterplattenlayout der Bridge. Der linke Bereich um die XLR-Buchsen
        trägt keine Massefläche des Hauptteils; beide Massepotenziale bleiben
        getrennt (Abschnitt~\ref{subsec:bridge-pcb}).}
    \label{fig:pcb-bridge}
\end{figure}

\begin{figure}[htbp]
    \centering
    \begin{subfigure}[b]{0.32\linewidth}
        \centering
        \includegraphics[height=6.5cm]{images/anhang/pcb_tube_top.pdf}
        \caption{Tube, Oberseite}
    \end{subfigure}
    \hfill
    \begin{subfigure}[b]{0.32\linewidth}
        \centering
        \includegraphics[height=6.5cm]{images/anhang/pcb_tube_bot.pdf}
        \caption{Tube, Unterseite}
    \end{subfigure}
    \hfill
    \begin{subfigure}[b]{0.32\linewidth}
        \centering
        \includegraphics[height=6.5cm]{images/anhang/pcb_interface_top.pdf}
        \caption{Interface, Oberseite}
    \end{subfigure}
    \caption{Leiterplattenlayouts von Tube- und Interface-PCB. Der schmale
        Formfaktor folgt der Innengeometrie des Aluminiumprofils
        (Abschnitt~\ref{sec:gehaeuse}).}
    \label{fig:pcb-tube-interface}
\end{figure}

\begin{figure}[htbp]
    \centering
    \includegraphics[height=6.5cm]{images/anhang/pcb_interface_bot.pdf}
    \caption{Interface-PCB, Unterseite. Auf dieser Seite sind Display und Taster
        bestückt.}
    \label{fig:pcb-interface-bot}
\end{figure}

\clearpage

\section{Vollständige Messwerttabellen}
\label{sec:anhang-messwerte}

Kapitel~\ref{chap:evaluierung} druckt je Messreihe eine Auswahl der Kennzahlen
ab. Die folgenden Tabellen führen sämtliche berechneten Größen auf. Sie werden
aus denselben unveränderten Rohdaten erzeugt wie die Tabellen des Hauptteils
(Abschnitt~\ref{subsec:messgrenzen}).

\begin{table}[H]
    \centering
    \small
    \caption{Zuordnung der Rohdatendateien zu den Messpunkten. Alle Läufe wurden
        mit derselben Mess-Firmware auf WLAN-Kanal~1 über 300\,s aufgezeichnet
        (Abschnitt~\ref{subsec:messaufbau}).}
    \label{tab:messlaeufe}
    \begin{tabular}{llll}
        \toprule
        Rohdatendatei & Messreihe & Bedingung & Auswertung \\
        \midrule
        \texttt{range\_0m.csv}          & Reichweite & 0\,m, Sicht    & Abschnitt~\ref{sec:idealbedingungen}, \ref{sec:reichweitentest} \\
        \texttt{range\_10m.csv}         & Reichweite & 10\,m, Sicht   & Abschnitt~\ref{sec:idealbedingungen}, \ref{sec:reichweitentest} \\
        \texttt{range\_50m.csv}         & Reichweite & 50\,m, Sicht   & Abschnitt~\ref{sec:reichweitentest} \\
        \texttt{range\_100m.csv}        & Reichweite & 100\,m, Sicht  & Abschnitt~\ref{sec:reichweitentest} \\
        \texttt{range\_150m.csv}        & Reichweite & 150\,m, Sicht  & Abschnitt~\ref{sec:reichweitentest} \\
        \texttt{range\_200m.csv}        & Reichweite & 200\,m, Sicht  & Abschnitt~\ref{sec:reichweitentest} \\
        \midrule
        \texttt{range\_hidden\_10m.csv}  & Reichweite & 10\,m, verdeckt  & Abschnitt~\ref{sec:reichweitentest} \\
        \texttt{range\_hidden\_50m.csv}  & Reichweite & 50\,m, verdeckt  & Abschnitt~\ref{sec:reichweitentest} \\
        \texttt{range\_hidden\_100m.csv} & Reichweite & 100\,m, verdeckt & Abschnitt~\ref{sec:reichweitentest} \\
        \texttt{range\_hidden\_150m.csv} & Reichweite & 150\,m, verdeckt & Abschnitt~\ref{sec:reichweitentest} \\
        \texttt{range\_hidden\_200m.csv} & Reichweite & 200\,m, verdeckt & Abschnitt~\ref{sec:reichweitentest} \\
        \midrule
        \texttt{interference\_noisy.csv} & Störfestigkeit & Stoßzeit  & Abschnitt~\ref{sec:stoerfestigkeit} \\
        \texttt{interference\_calm.csv}  & Störfestigkeit & Nebenzeit & Abschnitt~\ref{sec:stoerfestigkeit} \\
        \bottomrule
    \end{tabular}
\end{table}

% Requires \usepackage{booktabs}; otherwise replace \toprule etc. with \hline.
\begin{table}[htbp]
  \centering
  \small
  \begin{tabular}{llrrrrrr}
    \toprule
    Kennzahl & Einheit & 0 m & 10 m & 50 m & 100 m & 150 m & 200 m \\
    \midrule
    Messdauer & s & 299 & 299 & 299 & 299 & 299 & 299 \\
    Erwartete Sequenzen &  & 11992 & 11992 & 11992 & 11992 & 11991 & 11992 \\
    Empfangene Sequenzen &  & 11992 & 11992 & 11985 & 11981 & 11654 & 10720 \\
    Verlorene Sequenzen &  & 0 & 0 & 7 & 11 & 337 & 1272 \\
    Sequenzverlust, kumuliert & \% & 0 & 0 & 0,06 & 0,09 & 2,81 & 10,61 \\
    Fensterverlust, Mittel & \% & 0 & 0 & 0,06 & 0,09 & 2,79 & 10,58 \\
    Fensterverlust, p95 & \% & 0 & 0 & 0 & 0 & 12,22 & 31,72 \\
    Fensterverlust, Maximum & \% & 0 & 0 & 2,50 & 7,14 & 43,75 & 60,47 \\
    Verlustfreie Fenster & \% & 100 & 100 & 97,67 & 97 & 61 & 8,67 \\
    Fenster mit Ausfall ab 50\,ms & \% & 0 & 0 & 0 & 0,33 & 8,67 & 40,67 \\
    Längste Verlustserie & Sequenzen & 0 & 0 & 1 & 3 & 11 & 17 \\
    Längste Verlustserie & ms & 0 & 0 & 25 & 75 & 275 & 425 \\
    Jitter, Mittel & ms & 4,14 & 6,16 & 5,82 & 6,62 & 8,71 & 14,94 \\
    Kopien je Sequenz, Mittel &  & 5,00 & 4,99 & 4,62 & 4,43 & 3,39 & 2,38 \\
    Kopien je Sequenz, Minimum &  & 4,88 & 4,55 & 4,15 & 3,52 & 1,87 & 1,39 \\
    RSSI, Mittel & dBm & $-$81,64 & $-$82,39 & $-$81,50 & $-$82,85 & $-$79,27 & $-$63,44 \\
    RSSI, Minimum & dBm & $-$93 & $-$96 & $-$96 & $-$95 & $-$94 & $-$94 \\
    \bottomrule
\end{tabular}
\caption{Vollständige Kennzahlen der Reichweitenmessung (mit Sichtverbindung). Auswahl daraus in Tabelle~\ref{tab:range-summary}.}
  \label{tab:range-full}
\end{table}


% Requires \usepackage{booktabs}; otherwise replace \toprule etc. with \hline.
\begin{table}[htbp]
  \centering
  \small
  \begin{tabular}{llrrrrr}
    \toprule
    Kennzahl & Einheit & 10 m & 50 m & 100 m & 150 m & 200 m \\
    \midrule
    Messdauer & s & 299 & 299 & 299 & 299 & 299 \\
    Erwartete Sequenzen &  & 11992 & 11993 & 11993 & 11992 & 11990 \\
    Empfangene Sequenzen &  & 11934 & 11743 & 10691 & 8732 & 7639 \\
    Verlorene Sequenzen &  & 58 & 250 & 1302 & 3260 & 4351 \\
    Sequenzverlust, kumuliert & \% & 0,48 & 2,08 & 10,86 & 27,18 & 36,29 \\
    Fensterverlust, Mittel & \% & 0,45 & 2,07 & 10,81 & 27,15 & 36,19 \\
    Fensterverlust, p95 & \% & 0 & 10,08 & 36,64 & 60,80 & 63,67 \\
    Fensterverlust, Maximum & \% & 30,91 & 41,46 & 87,80 & 82,93 & 85,71 \\
    Verlustfreie Fenster & \% & 97,67 & 64,33 & 25,67 & 0,67 & 0 \\
    Fenster mit Ausfall ab 50\,ms & \% & 1,67 & 7,33 & 37 & 80 & 95,67 \\
    Längste Verlustserie & Sequenzen & 3 & 9 & 23 & 28 & 38 \\
    Längste Verlustserie & ms & 75 & 225 & 575 & 700 & 950 \\
    Jitter, Mittel & ms & 6,96 & 9,23 & 16,45 & 26,59 & 32,12 \\
    Kopien je Sequenz, Mittel &  & 4,80 & 3,57 & 2,72 & 1,87 & 1,61 \\
    Kopien je Sequenz, Minimum &  & 2,34 & 1,33 & 1,20 & 1,17 & 1 \\
    RSSI, Mittel & dBm & $-$81,23 & $-$77,85 & $-$82,97 & $-$48,70 & $-$77,39 \\
    RSSI, Minimum & dBm & $-$95 & $-$95 & $-$95 & $-$94 & $-$93 \\
    \bottomrule
\end{tabular}
\caption{Vollständige Kennzahlen der Reichweitenmessung (ohne Sichtverbindung). Auswahl daraus in Tabelle~\ref{tab:range-hidden-summary}.}
  \label{tab:range-hidden-full}
\end{table}


% Requires \usepackage{booktabs}; otherwise replace \toprule etc. with \hline.
\begin{table}[htbp]
  \centering
  \small
  \begin{tabular}{llrr}
    \toprule
    Kennzahl & Einheit & Nebenzeit & Stoßzeit \\
    \midrule
    Messdauer & s & 299 & 299 \\
    Erwartete Sequenzen &  & 11992 & 11991 \\
    Empfangene Sequenzen &  & 11991 & 11930 \\
    Verlorene Sequenzen &  & 1 & 61 \\
    Sequenzverlust, kumuliert & \% & 0,01 & 0,51 \\
    Fensterverlust, Mittel & \% & 0,01 & 0,45 \\
    Fensterverlust, p95 & \% & 0 & 2,50 \\
    Fensterverlust, Maximum & \% & 2,50 & 56,60 \\
    Verlustfreie Fenster & \% & 99,67 & 89,67 \\
    Fenster mit Ausfall ab 50\,ms & \% & 0 & 0,33 \\
    Längste Verlustserie & Sequenzen & 1 & 30 \\
    Längste Verlustserie & ms & 25 & 750 \\
    Jitter, Mittel & ms & 6,47 & 8,45 \\
    Kopien je Sequenz, Mittel &  & 4,57 & 3,83 \\
    Kopien je Sequenz, Minimum &  & 4,03 & 3 \\
    RSSI, Mittel & dBm & $-$83,72 & $-$84,02 \\
    RSSI, Minimum & dBm & $-$97 & $-$98 \\
    \bottomrule
\end{tabular}
\caption{Vollständige Kennzahlen der Messreihe Störfestigkeit. Auswahl daraus in Tabelle~\ref{tab:interference-summary}.}
  \label{tab:interference-full}
\end{table}


\clearpage

\section{Messwerte Leistungsaufnahme}
\label{sec:anhang-leistungsaufnahme}

\begin{figure}[H]
    \centering
    \begin{subfigure}[b]{\linewidth}
        \centering
        \includegraphics[width=0.3\linewidth]{images/anhang/power_consumption_power_supply_only.jpeg}
        \caption{Leistungsaufnahme des Netzteils ohne angeschlossene Last.}
        \label{fig:power-supply-only}
    \end{subfigure}
\end{figure}

\begin{figure}[H]
    \centering
    \begin{subfigure}[b]{\linewidth}
        \centering
        \includegraphics[width=0.3\linewidth]{images/anhang/power_consumption_tube_no_light.jpeg}
        \caption{Leistungsaufnahme der Leuchtröhre ohne angeschlossene LEDs.}
        \label{fig:power-tube-no-light}
    \end{subfigure}
\end{figure}

\begin{figure}[H]
    \centering
    \begin{subfigure}[b]{\linewidth}
        \centering
        \includegraphics[width=0.3\linewidth]{images/anhang/power_consumption_tube_fully_on.jpeg}
        \caption{Leistungsaufnahme der Leuchtröhre mit angeschlossenen LEDs.}
        \label{fig:power-tube-fully-on}
    \end{subfigure}
\end{figure}

\section{Reproduktion der Ergebnisse}
\label{sec:anhang-reproduktion}

Firmware, Mess-Firmware, Rohdaten, Auswertungsskript, Leiterplattenentwürfe
einschließlich Fertigungsdaten und Gehäusedateien sind unter
\url{https://github.com/nichtrami/led_tube} veröffentlicht. Die gezeigten
Messergebnisse wurden mit dem Stand zum Zeitpunkt der Abgabe erhoben.

\begin{table}[H]
    \centering
    \small
    \caption{Verzeichnisstruktur des Repositorys.}
    \label{tab:repo-struktur}
    \begin{tabular}{ll}
        \toprule
        Verzeichnis & Inhalt \\
        \midrule
        \texttt{src/}            & Anwendungsfirmware beider Gerätetypen \\
        \texttt{test\_programs/} & Mess-Firmware, Rohdaten und Auswertung \\
        \texttt{pcb/}            & KiCad-Projekte und Gerber-Daten der drei Baugruppen \\
        \texttt{housings/}       & Gehäuse als 3MF-Dateien für den 3D-Druck \\
        \texttt{documents/}      & Datenblätter der eingesetzten Bauteile \\
        \bottomrule
    \end{tabular}
\end{table}

Der Gerätetyp wird über die PlatformIO-Zielumgebung gewählt
(Abschnitt~\ref{sec:entwicklungsumgebung}), die Bibliotheksabhängigkeiten sind
in \texttt{platformio.ini} versioniert.

\begin{lstlisting}[language=bash, caption={Übersetzen und Flashen.}, captionpos=b]
pio run -e led_tube -t upload    # Leuchtroehre
pio run -e bridge   -t upload    # Bridge
\end{lstlisting}

Die Mess-Firmware bindet per \texttt{build\_src\_filter} den produktiven
\texttt{WirelessDmxController} unverändert ein
(Abschnitt~\ref{sec:testaufbau}). \texttt{capture.sh} zeichnet die serielle
Ausgabe unverändert als CSV auf; der Dateiname bestimmt Testtyp und
Messbedingung. \texttt{analyze.py} erzeugt daraus die Abbildungen und Tabellen
dieser Arbeit neu.

\begin{lstlisting}[language=bash, caption={Messreihe aufzeichnen und auswerten.}, captionpos=b]
pio run -e test_reliability_bridge -t upload
pio run -e test_reliability_tube   -t upload

cd test_programs/results
./capture.sh raw/range_50m.csv 300      # 300 s aufzeichnen

cd analysis
pip install -r requirements.txt
python3 analyze.py                      # erzeugt figures/ und tables/
\end{lstlisting}
